\begin{center}
	\large\bfseries PERNYATAAN KODE ETIK \\ PENGGUNAAN AI GENERATIF \\
	\vspace{-0.3em}
	{\normalfont\footnotesize\itshape Code of Conduct Statement: Generative AI or AI-Assisted Usage}
\end{center}

\thispagestyle{empty}
\phantomsection
\addcontentsline{toc}{chapter}{PERNYATAAN KODE ETIK PENGGUNAAN AI GENERATIF}

\vspace{1ex} % Jarak diperkecil

\noindent Saya yang bertanda tangan di bawah ini: \\
{\footnotesize\itshape I, the undersigned:}

\vspace{0.5ex}

% --- BAGIAN IDENTITAS (Compact & Full Width) ---
\noindent
\begingroup
\renewcommand{\arraystretch}{1.1}
\begin{tabularx}{\textwidth}{@{}p{4.5cm} p{0.2cm} X@{}} % @{} menghapus padding kiri-kanan agar full margin
	Nama Mahasiswa / NRP \newline {\scriptsize\itshape Full Name / Student ID} & : & \name{} / \nrp{} \\

	Program Studi \newline {\scriptsize\itshape Study Program} & : & S-1 Teknik Informatika \\

	Judul Tugas Akhir \newline {\scriptsize\itshape Final Project Title} & : & \tatitle{} \\
\end{tabularx}
\endgroup

\vspace{1.5ex}

\noindent dengan ini menyatakan bahwa pada Tugas Akhir dengan judul di atas tersebut: \\
{\footnotesize\itshape hereby declare that in the Final Project with the above title:}

\vspace{0.5ex}

% Tabel
\begingroup
\small
\noindent
\begin{NiceTabularX}{\textwidth}{|c|X|c|}[cell-space-limits=2pt]
	\hline

	% HEADER
	\rowcolor{gray!25}

	% Kolom 1: No
	\Block{2-1}{\textbf{No.}} &

	% Kolom 2: Pernyataan
	\Block{2-1}{\textbf{Pernyataan} \\ \textit{\scriptsize Statement}} &

	% Kolom 3: Centang
	\Block{2-1}{(\emo{check-mark-button})} \\

	\rowcolor{gray!25} & & \\
	\hline

	% ISI
	1 & Saya hanya menggunakan AI generatif sebagai alat bantu untuk memperbaiki tata bahasa. AI generatif tidak digunakan untuk membuat isi Tugas Akhir. \newline
	{\scriptsize\itshape I only used generative AI as a tool to improve the readability or language of the text in my Final Project. It was not used to generate a complete text of my work.}
	& \emo{check-mark-button} \\ \hline

	2 & Saya telah memeriksa dan/atau memperbaiki seluruh bagian dari Tugas Akhir saya yang dibantu oleh AI generatif agar sesuai dengan baku mutu penulisan karya ilmiah. \newline
	{\scriptsize\itshape I have reviewed and refined all aspects of my work that generative AI assists with, ensuring it adheres to the standards of academic writing.}
	& \emo{check-mark-button} \\ \hline

	3 & Saya tidak menggunakan AI generatif untuk pembuatan data primer, grafik dan/atau tabel pada Tugas Akhir saya. \newline
	{\scriptsize\itshape I did not use generative AI to generate primary data, figures, and/or tables in my work.}
	& \emo{check-mark-button} \\ \hline

	4 & Saya telah memberikan atribusi/pengakuan terhadap alat AI yang digunakan, secara rinci pada suatu bagian pada lampiran. \newline
	{\scriptsize\itshape I have acknowledged the use of generative AI in any part of the work in the specific appendix page.}
	& \emo{check-mark-button} \\ \hline

	5 & Saya memastikan tidak ada plagiarisme, termasuk hal yang berasal dari penggunaan AI generatif. \newline
	{\scriptsize\itshape I have ensured that there is no plagiarism issue in the work, including any parts generated by AI.}
	& \emo{check-mark-button} \\ \hline

\end{NiceTabularX}
\endgroup

% Gunakan vfill agar tanda tangan turun ke bawah, tapi secukupnya
\vfill

% TANDA TANGAN
\noindent
\begin{tabularx}{\textwidth}{X r}
	& Surabaya, \MONTH{} \the\year{} \\
	& Mahasiswa \\
	& \\
	& \\
	& \\
	& \\ % Baris kosong untuk tanda tangan
	& \textbf{\name{}} \\
	& NRP. \nrp{}
\end{tabularx}

\clearpage